\documentclass[11pt,letterpaper]{article}

\usepackage[utf8]{inputenc}
\usepackage[spanish]{babel}
\usepackage[margin=2.5cm]{geometry}
\usepackage{amsmath,amssymb,amsthm}
\usepackage{enumitem}
\usepackage{titlesec}
\usepackage{hyperref}
\usepackage{xcolor}

\definecolor{unicaucaverde}{RGB}{0,90,60}
\hypersetup{colorlinks=true,linkcolor=unicaucaverde,urlcolor=unicaucaverde}

\titleformat{\section}{\normalfont\Large\bfseries\color{unicaucaverde}}{\thesection}{1em}{}
\titleformat{\subsection}{\normalfont\large\bfseries}{\thesubsection}{1em}{}

\theoremstyle{definition}
\newtheorem{definicion}{Definición}[section]
\newtheorem{ejemplo}[definicion]{Ejemplo}
\newtheorem{propiedad}[definicion]{Propiedad}
\theoremstyle{plain}
\newtheorem{teorema}[definicion]{Teorema}

\newcommand{\R}{\mathbb{R}}
\newcommand{\vv}[1]{\mathbf{#1}}

\title{\textcolor{unicaucaverde}{\Large Notas de clase --- Semana 2}\\[4pt]
  \large Eliminación de Gauss-Jordan y sistemas homogéneos}
\author{Álgebra Lineal para Ingeniería --- Universidad del Cauca}
\date{2026}

\begin{document}
\maketitle
\tableofcontents

\section{Sistemas equivalentes y clasificación}

\begin{definicion}[Sistemas equivalentes]
Dos sistemas de ecuaciones lineales son \emph{equivalentes} si tienen exactamente el mismo
conjunto de soluciones. Que dos sistemas sean equivalentes \emph{no} significa que sean iguales:
pueden tener ecuaciones distintas y, sin embargo, toda solución de uno ser solución del otro.
\end{definicion}

\begin{teorema}[Las operaciones elementales de fila producen sistemas equivalentes]\label{teo:equiv}
Cada una de las tres operaciones elementales de fila aplicada a la matriz aumentada
$[A\mid\vv{b}]$ produce la matriz aumentada de un sistema equivalente al original:
\begin{enumerate}[label=(\roman*)]
  \item Intercambiar dos filas (equivale a reordenar ecuaciones).
  \item Multiplicar una fila por un escalar $c\neq 0$ (equivale a multiplicar ambos lados de una
        ecuación por $c$).
  \item Sumar a una fila un múltiplo escalar de otra (equivale a reemplazar una ecuación por su
        suma con un múltiplo de otra).
\end{enumerate}
Por tanto, el sistema obtenido al final del algoritmo de Gauss-Jordan tiene el mismo conjunto
solución que el sistema original.
\end{teorema}

\begin{proof}[Justificación]
Cada operación es \emph{reversible}: se puede deshacer con otra operación del mismo tipo. Como
cada paso preserva el conjunto solución y se puede volver al sistema anterior, los dos sistemas
tienen exactamente las mismas soluciones, es decir, son equivalentes.
\end{proof}

\begin{definicion}[Sistema consistente e inconsistente]
\begin{itemize}
  \item Un sistema es \emph{consistente} si tiene \textbf{al menos una solución} (puede ser única
        o infinitas).
  \item Un sistema es \emph{inconsistente} si \textbf{no tiene ninguna solución}.
\end{itemize}
La inconsistencia se detecta en la forma escalonada: existe una fila del tipo
$[\,0\;\cdots\;0\mid c\,]$ con $c\neq 0$, que representa la ecuación $0=c$, una contradicción.
\end{definicion}

\begin{ejemplo}
El sistema
\[
  x_1 + x_2 = 3, \qquad 2x_1 + 2x_2 = 7
\]
es inconsistente. Al reducir la matriz aumentada:
\[
  \left(\begin{array}{cc|c}1&1&3\\2&2&7\end{array}\right)
  \xrightarrow{F_2-2F_1}
  \left(\begin{array}{cc|c}1&1&3\\0&0&1\end{array}\right)
\]
La segunda fila dice $0=1$: contradicción. El sistema no tiene solución.
\end{ejemplo}

\begin{ejemplo}[Caso 1: sistema inconsistente]
\[
  x_1 + x_2 = 3, \qquad 2x_1 + 2x_2 = 7.
\]
Reducción de la matriz aumentada:
\[
  \left(\begin{array}{cc|c}1&1&3\\2&2&7\end{array}\right)
  \xrightarrow{F_2 - 2F_1}
  \left(\begin{array}{cc|c}1&1&3\\0&0&1\end{array}\right).
\]
La fila $[0\ 0\mid 1]$ representa la ecuación $0=1$: \textbf{contradicción}. El sistema es
\textbf{inconsistente}; no tiene solución. (Geométricamente, las dos rectas son paralelas.)
\end{ejemplo}

\begin{ejemplo}[Caso 2: solución única]
\[
  x_1 + x_2 + x_3 = 6,\quad 2x_1 - x_2 + x_3 = 3,\quad x_1 + 2x_2 - x_3 = 2.
\]
Reducción completa (RREF):
\[
  \left(\begin{array}{ccc|c}1&1&1&6\\2&-1&1&3\\1&2&-1&2\end{array}\right)
  \xrightarrow{F_2-2F_1,\;F_3-F_1}
  \left(\begin{array}{ccc|c}1&1&1&6\\0&-3&-1&-9\\0&1&-2&-4\end{array}\right)
\]
\[
  \xrightarrow{F_2\leftrightarrow F_3}
  \left(\begin{array}{ccc|c}1&1&1&6\\0&1&-2&-4\\0&-3&-1&-9\end{array}\right)
  \xrightarrow{F_3+3F_2}
  \left(\begin{array}{ccc|c}1&1&1&6\\0&1&-2&-4\\0&0&-7&-21\end{array}\right)
\]
\[
  \xrightarrow{-\tfrac{1}{7}F_3}
  \left(\begin{array}{ccc|c}1&1&1&6\\0&1&-2&-4\\0&0&1&3\end{array}\right)
  \xrightarrow{\substack{F_2+2F_3\\F_1-F_3}}
  \left(\begin{array}{ccc|c}1&1&0&3\\0&1&0&2\\0&0&1&3\end{array}\right)
  \xrightarrow{F_1-F_2}
  \left(\begin{array}{ccc|c}1&0&0&1\\0&1&0&2\\0&0&1&3\end{array}\right).
\]
Rango $r=3=n$: el sistema es \textbf{consistente con solución única}
\[
  \boxed{x_1=1,\quad x_2=2,\quad x_3=3.}
\]
\end{ejemplo}

\begin{ejemplo}[Caso 3: infinitas soluciones]
\[
  x_1 + 2x_2 - x_3 = 0, \qquad 2x_1 + 4x_2 - 2x_3 = 0.
\]
(La segunda ecuación es el doble de la primera.) Reducción:
\[
  \left(\begin{array}{ccc|c}1&2&-1&0\\2&4&-2&0\end{array}\right)
  \xrightarrow{F_2-2F_1}
  \left(\begin{array}{ccc|c}1&2&-1&0\\0&0&0&0\end{array}\right).
\]
Rango $r=1 < n=3$: hay $n-r=2$ variables libres. Tomando $x_2=s$ y $x_3=t$ libres:
\[
  x_1 = -2s + t.
\]
El conjunto solución (familia de infinitas soluciones) es
\[
  \boxed{\{(-2s+t,\;s,\;t)\mid s,t\in\mathbb{R}\}.}
\]
El vector $(0,0,0)$ corresponde a $s=t=0$; p. ej., $(1,0,1)$ se obtiene con $s=0,\,t=1$.
\end{ejemplo}


\section{\texorpdfstring{$m$}{m} ecuaciones con \texorpdfstring{$n$}{n} incógnitas}

\begin{definicion}[Matriz aumentada]
Al sistema de $m$ ecuaciones con $n$ incógnitas
\[
\begin{aligned}
  a_{11}x_1 + \cdots + a_{1n}x_n &= b_1 \\
  &\;\vdots \\
  a_{m1}x_1 + \cdots + a_{mn}x_n &= b_m
\end{aligned}
\]
se le asocia la matriz aumentada
\[
  [A \mid \vv{b}] =
  \left(\begin{array}{ccc|c}
    a_{11} & \cdots & a_{1n} & b_1 \\
    \vdots & \ddots & \vdots & \vdots \\
    a_{m1} & \cdots & a_{mn} & b_m
  \end{array}\right),
\]
que contiene toda la información del sistema sin escribir las incógnitas.
\end{definicion}

\begin{definicion}[Operaciones elementales de fila]
Sobre la matriz aumentada se permiten tres operaciones, cada una equivalente a una manipulación
algebraica válida del sistema (no cambian el conjunto solución):
\begin{enumerate}[label=(\roman*)]
  \item Intercambiar dos filas.
  \item Multiplicar una fila por un escalar distinto de cero.
  \item Sumar a una fila un múltiplo escalar de otra fila.
\end{enumerate}
\end{definicion}

\begin{definicion}[Forma escalonada reducida por filas]
Una matriz está en \emph{forma escalonada reducida} (RREF, por sus siglas en inglés) si:
\begin{enumerate}[label=(\roman*)]
  \item Toda fila no nula tiene como primer elemento no nulo (el \emph{pivote}) un 1.
  \item Cada pivote está a la derecha del pivote de la fila anterior.
  \item Las filas nulas (si las hay) están al final.
  \item En la columna de cada pivote, todos los demás elementos son 0.
\end{enumerate}
\end{definicion}

\begin{definicion}[Eliminación de Gauss-Jordan]
Es el algoritmo que, mediante una secuencia de operaciones elementales de fila, lleva la matriz
aumentada $[A\mid\vv{b}]$ a su forma escalonada reducida. Las incógnitas que corresponden a una
columna con pivote quedan determinadas directamente; las que no tienen pivote son
\emph{variables libres}, y el sistema tiene infinitas soluciones parametrizadas por ellas. La
\emph{eliminación gaussiana} es la variante que solo lleva la matriz a forma escalonada (no
reducida) y luego resuelve por sustitución hacia atrás --- mismo resultado, menos operaciones
pero un paso final adicional.
\end{definicion}

\begin{teorema}[Número de soluciones]
Sea $r$ el número de pivotes (filas no nulas) de la forma escalonada de $[A\mid\vv{b}]$, y sea
$r'$ el número de pivotes de $A$ sola. El sistema:
\begin{itemize}
  \item es \textbf{inconsistente} si $r' < r$ (hay un pivote en la columna de $\vv{b}$, es decir
        una fila del tipo $[0\;\cdots\;0\mid c]$ con $c\neq 0$);
  \item tiene \textbf{solución única} si es consistente y $r=n$ (no hay variables libres);
  \item tiene \textbf{infinitas soluciones} si es consistente y $r<n$ (hay $n-r$ variables
        libres).
\end{itemize}
\end{teorema}

\begin{ejemplo}
Resolver por Gauss-Jordan:
\[
\begin{aligned}
  x_1 + x_2 + x_3 &= 6 \\
  2x_1 - x_2 + x_3 &= 3 \\
  x_1 + 2x_2 - x_3 &= 2
\end{aligned}
\]
Matriz aumentada y reducción:
\[
  \left(\begin{array}{ccc|c} 1&1&1&6\\ 2&-1&1&3\\ 1&2&-1&2 \end{array}\right)
  \xrightarrow{F_2-2F_1,\; F_3-F_1}
  \left(\begin{array}{ccc|c} 1&1&1&6\\ 0&-3&-1&-9\\ 0&1&-2&-4 \end{array}\right)
\]
\[
  \xrightarrow{F_2\leftrightarrow F_3}
  \left(\begin{array}{ccc|c} 1&1&1&6\\ 0&1&-2&-4\\ 0&-3&-1&-9 \end{array}\right)
  \xrightarrow{F_3+3F_2}
  \left(\begin{array}{ccc|c} 1&1&1&6\\ 0&1&-2&-4\\ 0&0&-7&-21 \end{array}\right)
\]
De la última fila, $x_3=3$; retrocediendo (o continuando a RREF), $x_2=-4+2(3)=2$, y
$x_1=6-2-3=1$. Solución única: $(1,2,3)$.
\end{ejemplo}

\subsection{Ejemplo de aplicación: corrientes de malla en un circuito}\label{sec:circuito}

Un circuito de tres mallas tiene, en la malla 1, una resistencia propia de $2\,\Omega$ y una
fuente de $5\,\text{V}$; la malla 2 comparte $1\,\Omega$ con la malla 1 y $1\,\Omega$ con la
malla 3, más $1\,\Omega$ propio (resistencia total $3\,\Omega$), sin fuente; la malla 3 tiene
resistencia propia $2\,\Omega$ y una fuente de $1\,\text{V}$. Al aplicar la ley de voltajes de
Kirchhoff (LVK) a cada malla se obtiene el sistema, en las corrientes de malla
$I_1,I_2,I_3$ (amperios):
\[
\begin{aligned}
  2I_1 - I_2 &= 5 \\
  -I_1 + 3I_2 - I_3 &= 0 \\
  -I_2 + 2I_3 &= 1
\end{aligned}
\]

\paragraph{Solución sin un método sistemático (sustitución).}
De la primera ecuación, $I_2 = 2I_1-5$. De la tercera, $I_3 = \dfrac{1+I_2}{2} =
\dfrac{2I_1-4}{2} = I_1-2$. Sustituyendo ambas en la segunda ecuación:
\[
  -I_1 + 3(2I_1-5) - (I_1-2) = 0 \;\Longrightarrow\; 4I_1 - 13 = 0 \;\Longrightarrow\; I_1 = \tfrac{13}{4}.
\]
Con $I_1=\tfrac{13}{4}$: $I_2 = 2\cdot\tfrac{13}{4}-5 = \tfrac{3}{2}$, y
$I_3 = \tfrac{13}{4}-2 = \tfrac{5}{4}$. Este método funcionó, pero requirió decidir en qué orden
despejar cada variable; con más de tres incógnitas se vuelve difícil de organizar y propenso a
errores. La eliminación de Gauss-Jordan resuelve el mismo problema con un procedimiento fijo,
sin decisiones de por medio.

\paragraph{Solución por Gauss-Jordan (con fracciones exactas).}
\[
  \left(\begin{array}{ccc|c} 2&-1&0&5\\ -1&3&-1&0\\ 0&-1&2&1 \end{array}\right)
  \xrightarrow{\;\frac{1}{2}F_1\;}
  \left(\begin{array}{ccc|c} 1&-\tfrac{1}{2}&0&\tfrac{5}{2}\\ -1&3&-1&0\\ 0&-1&2&1 \end{array}\right)
  \xrightarrow{F_2+F_1}
  \left(\begin{array}{ccc|c} 1&-\tfrac{1}{2}&0&\tfrac{5}{2}\\ 0&\tfrac{5}{2}&-1&\tfrac{5}{2}\\ 0&-1&2&1 \end{array}\right)
\]
\[
  \xrightarrow{\frac{2}{5}F_2}
  \left(\begin{array}{ccc|c} 1&-\tfrac{1}{2}&0&\tfrac{5}{2}\\ 0&1&-\tfrac{2}{5}&1\\ 0&-1&2&1 \end{array}\right)
  \xrightarrow{\substack{F_1+\frac{1}{2}F_2 \\ F_3+F_2}}
  \left(\begin{array}{ccc|c} 1&0&-\tfrac{1}{5}&3\\ 0&1&-\tfrac{2}{5}&1\\ 0&0&\tfrac{8}{5}&2 \end{array}\right)
\]
\[
  \xrightarrow{\frac{5}{8}F_3}
  \left(\begin{array}{ccc|c} 1&0&-\tfrac{1}{5}&3\\ 0&1&-\tfrac{2}{5}&1\\ 0&0&1&\tfrac{5}{4} \end{array}\right)
  \xrightarrow{\substack{F_1+\frac{1}{5}F_3 \\ F_2+\frac{2}{5}F_3}}
  \left(\begin{array}{ccc|c} 1&0&0&\tfrac{13}{4}\\ 0&1&0&\tfrac{3}{2}\\ 0&0&1&\tfrac{5}{4} \end{array}\right)
\]
Misma solución que por sustitución: $I_1=\tfrac{13}{4}$, $I_2=\tfrac{3}{2}$, $I_3=\tfrac{5}{4}$
amperios. Nótese que ninguna fracción se aproximó a decimal en ningún paso: trabajar con
fracciones exactas evita el error de redondeo que se acumularía si se truncaran los decimales en
cada operación --- relevante en cualquier cálculo de ingeniería hecho a mano o verificado a mano.

\section{Sistemas homogéneos de ecuaciones}

\begin{definicion}[Sistema homogéneo]
Un sistema de ecuaciones lineales es \emph{homogéneo} si todos los términos independientes son
cero:
\[
\begin{aligned}
  a_{11}x_1 + \cdots + a_{1n}x_n &= 0 \\
  &\;\vdots \\
  a_{m1}x_1 + \cdots + a_{mn}x_n &= 0
\end{aligned}
\]
En notación matricial, $A\vv{x}=\vv{0}$.
\end{definicion}

\begin{propiedad}[Solución trivial]
Todo sistema homogéneo es consistente, porque $\vv{x}=\vv{0}=(0,\dots,0)$ siempre lo satisface
(cada ecuación queda $0=0$). Esta se llama la \emph{solución trivial}. La pregunta relevante no es
si hay solución, sino si hay \textbf{alguna otra} (\emph{solución no trivial}).
\end{propiedad}

\begin{teorema}[Cuándo hay soluciones no triviales]
Un sistema homogéneo de $m$ ecuaciones con $n$ incógnitas tiene soluciones no triviales si y solo
si, al llevar $A$ a forma escalonada, el número de pivotes $r$ es menor que $n$ (hay al menos una
variable libre). En particular, si $m<n$ (menos ecuaciones que incógnitas), el sistema homogéneo
\textbf{siempre} tiene soluciones no triviales.
\end{teorema}

\begin{ejemplo}
\[
\begin{aligned}
  x_1 + 2x_2 - x_3 &= 0 \\
  2x_1 + 4x_2 - 2x_3 &= 0
\end{aligned}
\]
La segunda ecuación es el doble de la primera (mismo hiperplano), así que $r=1 < n=3$: hay
$n-r=2$ variables libres. Tomando $x_2=s$, $x_3=t$ libres: $x_1=-2s+t$. El conjunto solución es
$\{(-2s+t,\,s,\,t) : s,t\in\R\}$, infinito, incluyendo la trivial ($s=t=0$).
\end{ejemplo}

\section{Soluciones de los ejercicios propuestos}

\begin{description}

\item[\textbf{Ejercicio 1.}]
Sistema $\{2x_1-x_2+x_3=4,\; x_1+3x_3=-1,\; -x_1+x_2+2x_3=0\}$.
\[
  \left(\begin{array}{ccc|c}2&-1&1&4\\1&0&3&-1\\-1&1&2&0\end{array}\right)
  \xrightarrow{F_1\leftrightarrow F_2}
  \left(\begin{array}{ccc|c}1&0&3&-1\\2&-1&1&4\\-1&1&2&0\end{array}\right)
\]
\[
  \xrightarrow{F_2-2F_1,\;F_3+F_1}
  \left(\begin{array}{ccc|c}1&0&3&-1\\0&-1&-5&6\\0&1&5&-1\end{array}\right)
  \xrightarrow{F_3+F_2}
  \left(\begin{array}{ccc|c}1&0&3&-1\\0&-1&-5&6\\0&0&0&5\end{array}\right)
\]
La fila $[0\;0\;0\mid 5]$ es la ecuación $0=5$: contradicción.
\[\boxed{\text{El sistema es \textbf{inconsistente}: no tiene solución.}}\]

\item[\textbf{Ejercicio 2.}]
Sistema $\{x+y+z=2,\;2x-y+z=1,\;x+2y-z=3\}$.
\[
  \left(\begin{array}{ccc|c}1&1&1&2\\2&-1&1&1\\1&2&-1&3\end{array}\right)
  \xrightarrow{F_2-2F_1,\;F_3-F_1}
  \left(\begin{array}{ccc|c}1&1&1&2\\0&-3&-1&-3\\0&1&-2&1\end{array}\right)
\]
\[
  \xrightarrow{F_2\leftrightarrow F_3}
  \left(\begin{array}{ccc|c}1&1&1&2\\0&1&-2&1\\0&-3&-1&-3\end{array}\right)
  \xrightarrow{F_3+3F_2}
  \left(\begin{array}{ccc|c}1&1&1&2\\0&1&-2&1\\0&0&-7&0\end{array}\right)
\]
\[
  \xrightarrow{-\frac{1}{7}F_3}
  \left(\begin{array}{ccc|c}1&1&1&2\\0&1&-2&1\\0&0&1&0\end{array}\right)
  \xrightarrow{\substack{F_2+2F_3\\F_1-F_3}}
  \left(\begin{array}{ccc|c}1&1&0&2\\0&1&0&1\\0&0&1&0\end{array}\right)
  \xrightarrow{F_1-F_2}
  \left(\begin{array}{ccc|c}1&0&0&1\\0&1&0&1\\0&0&1&0\end{array}\right)
\]
\[\boxed{x=1,\quad y=1,\quad z=0.}\]

\item[\textbf{Ejercicio 3.}]
Sistema $\{x+ky=1,\;kx+y=1\}$. La matriz aumentada reduce así:
\[
  \left(\begin{array}{cc|c}1&k&1\\k&1&1\end{array}\right)
  \xrightarrow{F_2-kF_1}
  \left(\begin{array}{cc|c}1&k&1\\0&1-k^2&1-k\end{array}\right)
\]
Nótese que $1-k^2=(1-k)(1+k)$ y $1-k=(1-k)\cdot 1$.
\begin{itemize}
  \item Si $k=1$: fila 2 queda $[0\;0\mid 0]$ →
        \textbf{infinitas soluciones}: $x+y=1$ (una fila de ceros).
  \item Si $k=-1$: fila 2 queda $[0\;0\mid 2]$ →
        \textbf{inconsistente}: $0=2$, sin solución.
  \item Si $k\neq\pm 1$: $1-k^2\neq 0$, se divide y se obtiene
        \[\boxed{x=y=\frac{1}{1+k}\quad(k\neq\pm 1).}\]
\end{itemize}

\item[\textbf{Ejercicio 4.}]
El sistema tiene $m=3$ ecuaciones y $n=5$ incógnitas ($m < n$). Por el teorema de soluciones no
triviales, todo sistema homogéneo con más incógnitas que ecuaciones tiene al menos una variable
libre, luego tiene soluciones no triviales. (El rango máximo posible es $r\leq m=3 < 5=n$, así que
siempre habrá $n-r\geq 2$ variables libres.)

\item[\textbf{Ejercicio 5.}]
Sistema homogéneo $\{x_1-x_2+2x_3=0,\;2x_1-2x_2+4x_3=0\}$ (la segunda es el doble de la
primera):
\[
  \left(\begin{array}{ccc|c}1&-1&2&0\\2&-2&4&0\end{array}\right)
  \xrightarrow{F_2-2F_1}
  \left(\begin{array}{ccc|c}1&-1&2&0\\0&0&0&0\end{array}\right)
\]
Rango $r=1$: \textbf{variables libres} $x_2=s$ y $x_3=t$.\\
De la primera fila: $x_1 = x_2 - 2x_3 = s - 2t$.
\[\boxed{\text{Conjunto solución: }\{(s-2t,\;s,\;t)\mid s,t\in\mathbb{R}\}.}\]
Ejemplos: $(0,0,0)$ (trivial, $s=t=0$);\; $(1,1,0)$ ($s=1, t=0$);\; $(-2,0,1)$ ($s=0,t=1$).

\item[\textbf{Ejercicio 6.}]
Un sistema homogéneo $3\times 3$ con solución trivial única no tiene variables libres. Como
$n=3$ incógnitas, el número de pivotes debe ser $r=3$ (ninguna columna sin pivote). La forma
escalonada de la matriz de coeficientes $A$ (sin la columna de ceros, que siempre queda igual)
tiene \textbf{3 pivotes}, uno por fila. (Equivalentemente: $A$ es no singular, $\det A\neq 0$.)

\end{description}

\end{document}
